%----------------------------------------------------------------------------------------
% БҮЛЭГ 5: ЗАГВАРЫН СУРГАЛТ БА ХЯНАЛТ
%----------------------------------------------------------------------------------------

\chapter{Загварын сургалт ба хяналт}
\label{Chapter_Training}

Энэ бүлэгт багш--сурагч загварын мэдлэг дамжуулалтын урсгалын бодит хэрэгжүүлэлт, сургалтын явц болон хяналтын дэлгэрэнгүй мэдээллийг тусгасан болно. Арга зүйн бүлэгт онолын үндэслэл, өгөгдлийн бүлэгт бэлтгэлийн ажил тусгагдсан бол энэ бүлэгт тэдгээрийг нэгтгэн ажиллуулсан бодит үйл явц, гарсан үр дүнгийн ажиглалт, гүйцэтгэлийн шинжилгээг дэлгэрэнгүй авч үзнэ.

%-------------------------------------------------------------------------------
\section{Сургалтын орчны бэлтгэл}
%-------------------------------------------------------------------------------

\subsection{Техник хангамж ба программ хангамж}

Бүх сургалтын туршилтыг NVIDIA A100 80GB SXM4 GPU дээр явуулсан. Тус GPU нь transformer загварын сургалтад тохиромжтой BF16 тооцооллыг техник хангамжийн түвшинд хурдасгадаг ба Flash Attention 2 дэмжлэгтэй.

\begin{table}[htbp]
\centering
\caption{Сургалтын орчны дэлгэрэнгүй тохиргоо}
\label{tab:train_env_full}
\begin{tabular}{ll}
\hline
\textbf{Бүрдэлхүүн} & \textbf{Тодорхойлолт} \\
\hline
GPU загвар              & NVIDIA A100 80GB SXM4 \\
GPU VRAM                & 80 GB HBM2e \\
CPU                     & Intel Xeon Gold 32 core \\
RAM                     & 128 GB DDR4 ECC \\
Хадгалалтын сан         & 2 TB NVMe SSD \\
CUDA хувилбар           & 11.8 \\
Python                  & 3.10.12 \\
PyTorch                 & 2.7.1+cu118 \\
Transformers            & 5.3.0 \\
TRL (Нарийвчилсан тохируулга)       & 0.12.0 \\
PEFT (LoRA)             & 0.14.0 \\
BitsAndBytes            & 0.45.0 \\
WandB (хяналт)          & 0.19.1 \\
\hline
\end{tabular}
\end{table}

\subsection{Орчин тохируулга ба хамаарлуудын суурилуулалт}

\begin{sloppypar}Сургалтын орчныг тохируулахдаа виртуал орчин ашиглан хамаарлуудын хувилбарын тогтвортой байдлыг хангасан. Бусад системтэй зөрчил үүсгэхгүйн тулд conda орчин ашигласан. Хамгийн чухал хамаарлуудыг дараах тэргүүлэх дарааллаар суурилуулсан: CUDA, PyTorch, transformers, peft, trl, bitsandbytes. BitsAndBytes хамаарлыг CUDA хэрэгслийн багцтай зохицуулан тусгайлан хөрвүүлсэн нь 4-bit квантчилалд шаардлагатай.\end{sloppypar}

Санамсаргүй тооны эхлэлийн утгыг бүх санд нэгэн зэрэг тогтоосон:
\begin{itemize}
    \item \texttt{torch.manual\_seed(42)}.
    \item \texttt{numpy.random.seed(42)}.
    \item \texttt{random.seed(42)}.
    \item \texttt{TRANSFORMERS\_SEED=42} хувьсагч.
\end{itemize}

Энэ нь туршилтыг ижил нөхцөлд давтан гүйцэтгэх боломжтой болгосон.

%-------------------------------------------------------------------------------
\section{Багш загварын дүгнэлт гаргалтын явц}
%-------------------------------------------------------------------------------

\subsection{Qwen3-VL-30B-A3B ачааллах ба тохируулах}

Багш загварыг ачааллахдаа \texttt{bfloat16} нарийвчлал болон \texttt{eager} attention mode-ийг ашигласан. Энэ тохиргоо нь \texttt{flash\_attention\_2}-тай харьцуулахад тогтвортой дүгнэлт гаргалт хийдэг боловч санах ойн хэрэглээ арай их байдаг. Загварыг ачааллах явцад GPU VRAM хэрэглээ дараах байдлаар өссөн: ачаалалт эхлэхэд 0 GB, загварын жин ачаалагдсаны дараа 54.3 GB, KV cache нэмэгдсэний дараа 62.2 GB.

\begin{figure}[H]
\centering
\includegraphics[width=0.9\textwidth]{Figures/model_inference.jpg}
\caption{Qwen3-VL-30B-A3B багш загварын дүгнэлт гаргалт явж байгаа байдал.}
\label{fig:teacher_inference_detail}
\end{figure}

\subsection{Дүгнэлт гаргалтын script (process\_v5\_ui.py)}

багш загварын дүгнэлт гаргалтын script нь дараах өвөрмөц онцлогуудтай:

\begin{itemize}
    \item \textbf{Rich terminal хяналтын самбар:} Боловсруулсан видео тоо болон GPU VRAM хэрэглээг нэгдсэн байдлаар харуулна.

    \item \textbf{Үргэлжлүүлэх механизм:} Тасалдсан тохиолдолд \texttt{train.jsonl} файлд хамгийн сүүлд бичигдсэн \texttt{source\_file}-ийг уншиж, тэр цэгээс үргэлжлүүлнэ. Энэ нь 34 цагийн урт дүгнэлт гаргалтад маш чухал байсан.

    \item \textbf{WandB интеграц:} \texttt{police-report-finetune} project дор бүх дүгнэлт гаргалтын метрикүүдийг бүртгэнэ.

    \item \textbf{Алдаа боловсруулалт:} Нэг видеоны дүгнэлт гаргалтад алдаа гарсан тохиолдолд тухайн файлыг алгасан, дараагийн файлд шилжих алдаа зохицуулалт.
\end{itemize}

%-------------------------------------------------------------------------------
\section{WandB тасралтгүй хяналт ба хяналтын систем}
%-------------------------------------------------------------------------------

\subsection{WandB төслийн тохиргоо}

Weights \& Biases (WandB) платформ дээр \texttt{police-report-finetune} нэртэй тусдаа төсөл үүсгэн бүх туршилтуудыг нэгдсэн системд бүртгэсэн. WandB-д дараах метрикүүдийг автоматаар бүртгэсэн:

\begin{itemize}
    \item \textbf{Сургалтын метрикийн бүртгэл:} train/loss, train/learning\_rate, train/epoch,\\ train/global\_step (10 алхам бүрт).
    \item \textbf{Eval метрикийн бүртгэл:} eval/loss, eval/runtime,\\ eval/samples\_per\_second (500 алхам бүрт).
    \item \textbf{Систем метрикийн бүртгэл:} GPU VRAM хэрэглээ, GPU ашиглалтын хувь,\\ CPU ашиглалт, RAM хэрэглээ (60 секунд бүрт).
    \item \textbf{Градиент нормын бүртгэл:} LoRA тохируулагчийн gradient norm\\ (хэт тааруулалт/exploding gradient илрүүлэхэд).
\end{itemize}

%-------------------------------------------------------------------------------
\section{Хадгалалтын цэг стратеги ба сургалт дахин эхлүүлэлт}
%-------------------------------------------------------------------------------

\subsection{Хадгалалтын цэг хадгалах стратеги}

500 алхам бүрт хадгалалтын цэг хадгалах дараах стратегийг баримталсан:

\begin{itemize}
    \item \textbf{Сүүлийн 10 хадгалалтын цэг хадгалах:} \texttt{save\_total\_limit=10} тохиргоог ашиглан хуучин хадгалалтын цэгүүдийг автоматаар устгадаг. Энэ нь дискний зай хэмнэхийн зэрэгцээ сүүлийн 5,000 алхмын хадгалалтын цэгүүдэд буцах боломжийг олгодог.
    \item \textbf{Eval loss-д суурилсан best model:} \texttt{load\_best\_model\_at\_end=True} тохиргоотой тул сургалт дуусахад хамгийн бага eval loss-тэй хадгалалтын цэг автоматаар ачаалагдана.
    \item \textbf{HuggingFace Hub синхрончилалт:} Сургалт дуусаад LoRA adapter-ийг HuggingFace Hub-т байршуулсан.
    \newline{\small\href{https://huggingface.co/stillmeta/fight-video-qlora1}{\texttt{fight-video-qlora1}}}
\end{itemize}

\subsection{Сургалт тасалдсан тохиолдолд дахин эхлүүлэх}

Урт сургалтын явцад тасалдал гарах эрсдэлтэй тул resume механизм тохируулсан:

\begin{enumerate}
    \item \texttt{trainer.train(resume\_from\_checkpoint=True)} параметр ашиглан хамгийн сүүлийн хадгалалтын цэгээс сургалт үргэлжлүүлэх.
    \item Optimizer state болон scheduler state-ийг хадгалалтын цэгээс ачааллах.
    \item Random state-ийг сэргээн тогтворт үр дүн хангах.
\end{enumerate}

Энэ механизм нь бодит сургалтын явцад гэнэтийн тасалдал гарсан тохиолдолд нэн чухал байсан.

%-------------------------------------------------------------------------------
\section{Загварын нэгтгэл ба байршуулалт}
%-------------------------------------------------------------------------------

\subsection{LoRA adapter нэгтгэх}

Сургалт дуусаад LoRA adapter-ийн жинг суурь загвартай нэгтгэх (merge) үйл явц хийгдэнэ. Нэгтгэлийн математик:

\begin{equation}
    W_{\text{merged}} = W_0 + \frac{\alpha}{r} \cdot B \cdot A
\end{equation}

Нэгтгэлийн дараа загвар нь стандарт Qwen2.5-VL-7B архитектуртай адил болж, тусгайлсан PEFT library байхгүйгээр ачааллах боломжтой. Нэгтгэгдсэн загварын хэмжээ нь суурь загварынхтай ижил ($\approx$14 GB BF16).

\subsection{Загварын байршуулалт ба ашиглалт}

Нэгтгэгдсэн загварыг дараах хоёр хэлбэрт байршуулсан:

\begin{itemize}
    \item \textbf{HuggingFace Hub байршуулалт:} LoRA adapter нь Hub-д байршуулагдсан.
    \newline{\small\href{https://huggingface.co/stillmeta/fight-video-qlora1}{\texttt{fight-video-qlora1}}}
    Энэ нь хамтрагчдад хялбарчлан хандах, хувилбарын хяналт хийх боломжийг олгодог.

    \item \textbf{Локал байршуулалт:} Нэгтгэгдсэн загвар нь локал SSD-д хадгалагдаж,\\ дүгнэлт гаргалт script-ийн шууд оролт болдог.
\end{itemize}

%-------------------------------------------------------------------------------
\section{Сургалтын явцад гарсан техникийн асуудал ба шийдэл}
%-------------------------------------------------------------------------------

Сургалтын явцад гарсан гол техникийн асуудлууд болон тэдгээрийн шийдлүүдийг дараах байдлаар нэгтгэлээ:

\begin{table}[H]
\centering
\caption{Сургалтын явцад тулгарсан асуудал ба шийдэл}
\label{tab:training_issues}
\small
\begin{tabularx}{\textwidth}{>{\raggedright\arraybackslash}p{3.2cm}>{\raggedright\arraybackslash}p{4.0cm}>{\raggedright\arraybackslash}X}
\hline
\textbf{Асуудал} & \textbf{Шалтгаан} & \textbf{Шийдэл} \\
\hline
Validation loss хэлбэлзэл & Eval batch size бага & Eval batch size 2 болгосон \\
Хадгалалтын цэг алдагдал & Серверийн хугацааны хязгаарлалт & Resume механизм нэмж, save\_steps 500 болгосон \\
Токены дараалал хэт урт & Зарим Q\&A хос 3,000+ токен & Max seq length 2,048 тогтоосон; урт Q\&A хосыг truncate хийсэн \\
\hline
\end{tabularx}
\end{table}

Эдгээр техникийн асуудлуудыг системтэйгээр шийдсэний үр дүнд 5 үеийн сургалт тогтвортой явагдаж, хүлээгдсэн үр дүн гарсан.
